% !TEX program = xelatex
\documentclass[10pt,a4paper]{article}

% ---- Пакеты для XeLaTeX и кириллицы ----
\usepackage{fontspec}
\usepackage{polyglossia}
\setmainlanguage{russian}
\setotherlanguage{english}

% ---- Шрифт с кириллицей ----
\setmainfont{Liberation Serif}
\setsansfont{Liberation Sans}
\setmonofont{Liberation Mono}

% ---- Поля и геометрия ----
\usepackage[top=1.5cm, bottom=1.5cm, left=1.5cm, right=1.5cm]{geometry}
\usepackage{parskip}
\setlength{\parindent}{0pt}

% ---- Продвинутые ссылки ----
\usepackage{hyperref}
\hypersetup{
    colorlinks=true,
    linkcolor=blue,
    urlcolor=blue,
    citecolor=blue
}

% ---- Полезные пакеты ----
\usepackage{enumitem}
\setlist{nosep, left=0pt}
\usepackage{tabularx}
\usepackage{array}
\usepackage{booktabs}
\usepackage{xcolor}
\usepackage{titlesec}

% ---- Настройка заголовков разделов ----
\titleformat{\section}{\Large\bfseries\uppercase}{}{0em}{}[\vspace{-0.3em}\rule{\textwidth}{0.5pt}]
\titlespacing*{\section}{0pt}{0.8em}{0.3em}

\titleformat{\subsection}{\large\bfseries}{}{0em}{}
\titlespacing*{\subsection}{0pt}{0.4em}{0.2em}

% ---- Пользовательские команды ----
\newcommand{\projectEntry}[3]{\textbf{#1} \hfill \textit{#2}\par #3\par\vspace{0.3em}}

% ---- Общие данные (имя, телефон, email и т.д.) ----
% ============================================================
% ОБЩИЕ ПЕРСОНАЛЬНЫЕ ДАННЫЕ
% ============================================================
% Меняешь здесь — меняется во всех резюме

\newcommand{\PersonName}{Кузык Иван Романович}
\newcommand{\PersonPhone}{+7 (916) 574-52-99}
\newcommand{\PersonEmail}{vanya.kuzyk1@gmail.com}
\newcommand{\PersonGitHub}{github.com/vanya890}
\newcommand{\PersonCity}{г.~Москва, м.~Бульвар Дмитрия Донского}
\newcommand{\PersonCitizenship}{РФ}
\newcommand{\PersonDrivingLicense}{Категория B}

% ============================================================
% КОМАНДА ДЛЯ ЗАГОЛОВКА РЕЗЮМЕ
% ============================================================
% Использование: \makeHeader{Должность}
\newcommand{\makeHeader}[1]{%
\begin{center}
    {\Huge\bfseries \PersonName}\\[0.3em]
    {\large #1}\\[0.5em]
    \begin{tabular}{ll}
        Телефон: & \PersonPhone \\
        Email:   & \href{mailto:\PersonEmail}{\PersonEmail} \\
        GitHub:  & \href{https://\PersonGitHub}{\PersonGitHub} \\
        Город:   & \PersonCity \\
        Гражданство: & \PersonCitizenship
    \end{tabular}
\end{center}
}

% ============================================================
% ОБЩЕЕ ОБРАЗОВАНИЕ
% ============================================================
% Использование: \makeEducation
\newcommand{\makeEducation}{%
\section{Образование}

\projectEntry{Московский политехнический университет --- Факультет информационных технологий}{2022--2026}
{Бакалавриат, направление «Информатика и вычислительная техника», специализация «Киберфизические системы». ВКР «Разработка образовательного 3D-симулятора робототехники» --- отлично. Средний балл --- 5.0.}

\projectEntry{ГБПОУ «Образовательный комплекс «Юго-Запад»}{2018--2022}
{Среднее специальное образование, специальность «Информационные системы и программирование». Диплом с отличием. Демонстрационный экзамен по компетенции «Разработка компьютерных игр и мультимедийных приложений» --- отлично.}
}

% ============================================================
% ОБЩАЯ ДОПОЛНИТЕЛЬНАЯ ИНФОРМАЦИЯ
% ============================================================
% Использование: \makeFooter{Ключевые компетенции}{Личные качества}
\newcommand{\makeFooter}[2]{%
\section{Дополнительная информация}

\begin{itemize}
    \item \textbf{GitHub:} \href{https://\PersonGitHub}{\PersonGitHub} --- #1
    \item \textbf{Ключевые компетенции:} #2
    \item \textbf{Личные качества:} Обучаемость, ответственность, умение работать с большим объёмом информации, системное мышление, ориентация на результат
    \item \textbf{Водительские права:} \PersonDrivingLicense
\end{itemize}
}

\begin{document}

% ============================================================
% ЗАГОЛОВОК
% ============================================================
\makeHeader{Инженер-техник / Специалист по CAD и робототехнике}

% ============================================================
% О СЕБЕ
% ============================================================
\section{О себе}

Технический специалист с образованием в области информационных систем и программирования. Имею опыт работы с инженерными CAD-системами (\textbf{Компас-3D, AutoCAD}), робототехникой (ROS, сервоприводы, лидары, дальномеры), системами технического зрения (OpenCV), VR/AR, физическими симуляциями (LabVIEW) и симуляторами промышленных ячеек. Работал с учётными системами (\textbf{1С, МойСклад}), занимался администрированием и переносом серверного оборудования. Выполнял проекты в рамках учёбы, практики и личных задач.

% ============================================================
% ОПЫТ РАБОТЫ (ПРОЕКТНАЯ ДЕЯТЕЛЬНОСТЬ)
% ============================================================
\section{Опыт работы}

\subsection{CAD-системы и чертежи}

\projectEntry{Работа с Компас-3D}{2023–2025}
{Создавал 3D-модели деталей и сборочных единиц в \textbf{Компас-3D} в рамках учебных проектов. Оформлял конструкторскую документацию и чертежи по ГОСТ.}

\projectEntry{Работа с AutoCAD}{2023–2024}
{Выполнял 2D-моделирование и создание планировок в \textbf{AutoCAD}. Оформлял чертежи в соответствии с требованиями.}

\projectEntry{Моделирование микросхем и дорожек}{2023–2024}
{Выполнял разводку печатных плат и работал с принципиальными схемами в рамках учебных проектов.}

\subsection{Робототехника и техническое зрение}

\projectEntry{Лабораторные работы с ROS (Clover \& TurtleBot3)}{2024–2025}
{Выполнил цикл лабораторных работ в среде \textbf{ROS (Robot Operating System)}:
\begin{itemize}
    \item Разработал полётное задание для дрона \textbf{Clover} с использованием маркеров ArUco
    \item Реализовал алгоритмы обхода препятствий с использованием \textbf{лидара}
    \item Запускал симуляцию \textbf{TurtleBot3} в Gazebo, работал с RViz
\end{itemize}}

\projectEntry{Система визуальной локализации робота (OpenCV)}{2024}
{Разработал desktop-приложение на \textbf{C\# Windows Forms} с использованием \textbf{OpenCvSharp} для определения положения и траектории робота по стереопотоку с двух камер. Реализовал захват кадров, обработку изображений, расчёт траектории.}

\projectEntry{Сборка и программирование роботов}{2023–2024}
{Работал с \textbf{сервоприводами}, датчиками линии, лидарами, дальномерами. Собирал и программировал простых роботов в рамках учебных проектов.}

\projectEntry{Симуляции промышленных ячеек}{2024}
{Работал с симуляторами \textbf{роботов-манипуляторов}, программировал траектории движения в рамках учебных проектов.}

\projectEntry{Физические симуляции (LabVIEW)}{2023–2024}
{Создавал виртуальные приборы и моделировал физические процессы в среде \textbf{LabVIEW}.}

\subsection{VR/AR}

\projectEntry{Разработка AR-порталов}{2024–2025}
{Создавал AR-порталы на базе маркеров и отслеживания движений. Работал с \textbf{веб-AR} через камеру браузера (JavaScript/Three.js/AR.js).}

\projectEntry{Системы виртуальной реальности}{2024}
{Работал с VR-очками, настраивал системы синхронного воспроизведения на нескольких VR-устройствах.}

\subsection{1С, учётные системы и администрирование}

\projectEntry{Оператор ПК (ГБПОУ «Юго-Запад»)}{Ноябрь 2020 – Январь 2021}
{Ввод и редактирование данных в программе \textbf{1С:Библиотека}, обучение сотрудников работе с системой на базовом уровне.}

\projectEntry{Работа с 1С:Торговая точка}{2023–2024}
{Настраивал кассовые места в \textbf{1С:Торговая точка}, работал с номенклатурой, базовое администрирование системы.}

\projectEntry{Работа с МойСклад}{2024}
{Работал с сервисом \textbf{МойСклад} для ведения складского учёта, настройка интеграций.}

\subsection{Сетевая инфраструктура}

\projectEntry{Перенос серверного оборудования (Гос. учреждение)}{2024}
{Участвовал в переносе серверного оборудования, занимался коммутацией и базовым администрированием сетевой инфраструктуры.}

\subsection{Практика}

\projectEntry{Техник-программист (ГБПОУ «Юго-Запад»)}{Апрель–Июнь 2022}
{Проверка работоспособности ПК, обновление и установка ОС и ПО, настройка сетевых подключений, установка принтеров и периферийного оборудования, помощь в организации ЕГЭ.}

\projectEntry{Оператор БД (Издательство «Просвещение»)}{Май–Июль 2021}
{Заполнение базы данных информацией, описание изображений для поиска по ключевым словам, подготовка данных для редакций.}

% ============================================================
% ОБРАЗОВАНИЕ
% ============================================================
\makeEducation

% ============================================================
% НАВЫКИ
% ============================================================
\section{Навыки}

\begin{tabular}{p{3.5cm} p{12cm}}
    \textbf{CAD/Чертежи:} & Компас-3D, AutoCAD, оформление конструкторской документации \\
    \textbf{Робототехника:} & ROS, Gazebo, RViz, Clover, TurtleBot3, сервоприводы, лидары, дальномеры, датчики линии \\
    \textbf{Техническое зрение:} & OpenCV (OpenCvSharp), обработка изображений, детекция объектов, трекинг \\
    \textbf{Симуляции:} & LabVIEW, симуляторы промышленных ячеек, физическое моделирование \\
    \textbf{VR/AR:} & Three.js/AR.js, веб-AR, маркерная AR, VR-очки \\
    \textbf{1С и учёт:} & 1С:Библиотека, 1С:Торговая точка, МойСклад \\
    \textbf{Сети:} & Коммутация, серверное оборудование, сетевая инфраструктура \\
    \textbf{Языки:} & C\#, Python, JavaScript, SQL \\
    \textbf{Инструменты:} & Git, Docker, Linux \\
    \textbf{Языки:} & Русский — родной, Английский — A2 (элементарный)
\end{tabular}

% ============================================================
% ДОПОЛНИТЕЛЬНАЯ ИНФОРМАЦИЯ
% ============================================================
\makeFooter{проекты по робототехнике, компьютерному зрению, симуляциям}{CAD-моделирование, робототехника, техническое зрение, VR/AR, администрирование 1С, сетевая инфраструктура}

\end{document}